\section{Formulae}

This section collects and derives all physical equations used in the simulation.

\subsection{Planetary Position}

The position of any body in the simulation is represented as a two-component vector:

\begin{equation}
\vec{r} = (x,\; y)
\end{equation}

\subsection{Distance Between Two Bodies}

The scalar separation $|\vec{r}|$ between the star at position $\vec{r}_\star$ and the planet at position $\vec{r}_p$ is given by the Euclidean norm of their displacement:

\begin{equation}
|\vec{r}| = \sqrt{(x_p - x_\star)^2 + (y_p - y_\star)^2}
\end{equation}

This is the standard Pythagorean formula extended to vector subtraction, and reduces to $\sqrt{x^2 + y^2}$ when the star is fixed at the origin.

\subsection{Two-Body Gravitational Acceleration}

The gravitational acceleration experienced by the planet due to the star is derived from Newton's Law of Universal Gravitation and his Second Law of Motion.

Newton's Law of Universal Gravitation states:

\begin{equation}
\vec{F} = G\,\frac{m_1\,m_2}{|\vec{r}|^3}\,\vec{r}
\end{equation}

where $\vec{r} = \vec{r}_p - \vec{r}_\star$ is the displacement vector from the star to the planet. Newton's Second Law gives:

\begin{equation}
\vec{F} = m\,\vec{a}
\end{equation}

Equating these and letting $m_1 = M$ (star mass) and $m_2 = m$ (planet mass):

\begin{align}
m\,\vec{a} &= G\,\frac{M\,m}{|\vec{r}|^3}\,\vec{r} \\[6pt]
\vec{a} &= G\,\frac{M}{|\vec{r}|^3}\,\vec{r}
\end{align}

The planet mass $m$ cancels exactly, confirming that the orbital dynamics are independent of the orbiting body's mass. Since gravity is attractive — the force acts \emph{toward} the star, opposite to the direction of $\vec{r}$ — a negative sign is introduced:

\begin{equation}
\boxed{\vec{a} = -\frac{GM}{|\vec{r}|^3}\,\vec{r}}
\end{equation}

With $G = 1$, this simplifies to $\vec{a} = -\dfrac{M}{|\vec{r}|^3}\,\vec{r}$.

\subsection{Specific Orbital Energy}

The specific orbital energy $\varepsilon$ (energy per unit mass of the planet) determines the type of trajectory:

\begin{equation}
\varepsilon = \frac{v^2}{2} - \frac{GM}{r}
\end{equation}

With $G = 1$ and using $M$ for the star's mass:

\begin{equation}
\varepsilon = \frac{v^2}{2} - \frac{M}{r}
\end{equation}

The sign of $\varepsilon$ classifies the orbit:
\begin{itemize}
    \item $\varepsilon < 0$: bound orbit (elliptical or circular)
    \item $\varepsilon = 0$: parabolic trajectory (marginally unbound)
    \item $\varepsilon > 0$: hyperbolic escape
\end{itemize}

\subsection{Orbital Velocities}

\subsubsection{Circular Orbital Velocity}

A perfectly circular orbit occurs when the gravitational centripetal acceleration exactly balances the required centripetal acceleration $v^2/r$:

\begin{equation}
\frac{v_c^2}{r} = \frac{GM}{r^2} \implies v_c = \sqrt{\frac{GM}{r}}
\end{equation}

\subsubsection{Escape Velocity}

The escape velocity is the minimum speed required for the planet to escape the star's gravitational field, i.e. the speed at which $\varepsilon = 0$:

\begin{equation}
\frac{v_e^2}{2} = \frac{GM}{r} \implies v_e = \sqrt{\frac{2GM}{r}}
\end{equation}

Note that $v_e = \sqrt{2}\,v_c$: escape velocity is always $\sqrt{2}$ times the circular velocity at the same radius.